% SGA 3, Expose VI_B, Loop 1 English reconstruction.
% Scope: Theorem 6.5.2 and Corollary 6.5.3 in full, including
% their complete proofs and source editor's notes 62--64.
% Authority: Polo--Gille Exp6B-13oct24.pdf, component-local pages
% 46--47 (combined-reader pages 387--388; displayed re-edition
% headers 376--377). Corollary 6.5.4 begins on local page 47 and
% is excluded.
% OCR and Reinhold English Markdown are comparison witnesses only.
% This unit contains no diagrams.

\medskip
\phantomsection\label{thm:sga3-VIB-6.5.2}%
\noindent\textbf{定理 6.5.2.---}
\emph{设 \(X\) 为 \(S\) 上的平滑概形, 其纤维几何
不可约.}%
\footnote{\label{note:sga3-VIB-62}编者注: 一方面,
我们将原文的 ``连通'' 改为
``不可约'' (见证明). 另一方面, 根据
[RG71], I, 3.3.4\textup{(iii)} 及 4.1.1, 只需假设
\(X\) 在 \(S\) 上平坦, 且纤维几何不可约
且无嵌入分支.}
\begin{enumerate}
\item[\textup{(i)}]\phantomsection\label{concl:sga3-VIB-6.5.2-i}
  \emph{对 \(X\) 的每个闭子概形 \(Y\), 函子
  \[
    \prod_{X/S}Y/X
  \]
  可由 \(S\) 的闭子概形 \(T\) 表出.}
\item[\textup{(ii)}]\phantomsection\label{concl:sga3-VIB-6.5.2-ii}
  \emph{此外, 若 \(Y\to X\) 有限呈示, 则
  \(T\to S\) 也有限呈示.}
\end{enumerate}

\emph{证明.}
用 \(f:X\to S\) 表示结构态射. 由于 \(f\)
忠实平坦且局部有限
呈示, 它是 fpqc 拓扑下的一个覆盖. 此外,
\[
  T=\prod_{X/S}Y/X
\]
显然是 fpqc 拓扑下 \(S\) 的子层. 因此,
\(T\) 能否由 \(S\) 的闭子概形表出,
是 \(S\) 上关于 fpqc 拓扑的局部问题,%
\footnote{\label{note:sga3-VIB-63}编者注: 例如见
第 VIII 篇增补的 1.7.}
而 \(T\) 在 \(S\) 上是否有限呈示的问题
也是如此 (EGA~IV\(_2\), 2.7.1).

作基变换 \(S'\to S\), 其中 \(S'=X\), 就化归为
\(X\) 有截面
\[
  \epsilon:S\longrightarrow X.
\]
的情形. 还可以假设 \(S\) 仿射, 因而
拟紧. 此时有:

\medskip
\phantomsection\label{cor:sga3-VIB-6.5.3}%
\noindent\textbf{推论 6.5.3.---}
\emph{在 \hyperref[thm:sga3-VIB-6.5.2]{6.5.2} 的假设下,
假设 \(S\) 拟紧, \(X\to S\) 有截面 \(\epsilon\),
且 \(Y\to X\) 有限呈示.
则存在一个整数 \(n\geq0\), 使得}
\phantomsection\label{eq:sga3-VIB-6.5.3-infinitesimal}%
\[
  \prod_{X/S}Y/X
  =
  \prod_{X_n/S}Y_n/X_n,
\]
\emph{其中 \(X_n\) 为浸入 \(\epsilon:S\to X\) 的第 \(n\) 阶
无穷小邻域, 且
\(Y_n=Y\cap X_n\).}

当 \(Y\) 在 \(X\) 上有限呈示时, 此推论
确实蕴涵 \hyperref[thm:sga3-VIB-6.5.2]{6.5.2}.
由于 \(X\) 在 \(S\) 上平滑, \(X_n\) 在 \(S\) 上有限且局部自由,
从而更是在 \(S\) 上本质自由
(见 \hyperref[def:sga3-VIB-6.2.1]{6.2.1}). 因此由
\hyperref[thm:sga3-VIB-6.2.3]{6.2.3},
\[
  \prod_{X_n/S}Y_n/X_n
\]
可由 \(S\) 的一个闭子概形 \(T\) 表出,
且它在 \(S\) 上有限呈示.

先在 \(S\) 为 Noether 概形时证明 \hyperref[cor:sga3-VIB-6.5.3]{6.5.3},
从而证明 \hyperref[thm:sga3-VIB-6.5.2]{6.5.2}. 令
\phantomsection\label{eq:sga3-VIB-6.5.3-Tn}%
\[
  T_n=\prod_{X_n/S}Y_n/X_n.
\]
\(T_n\) 构成 \(S\) 的一个递减闭子概形列,
由于 \(S\) 为 Noether 概形, 此列最终稳定.
设
\phantomsection\label{eq:sga3-VIB-6.5.3-R}%
\[
  R=\bigcap_{n\geq0}\prod_{X_n/S}Y_n/X_n
\]
为它们在 \(n\) 充分大时的共同值. 显然
\(T\subset R\), 只需证明 \(R\subset T\).

作基变换 \(R\to S\) 后, 化归为
\[
  R=S,
\]
即
\[
  Y_n=X_n\quad\text{对每个 \(n\)},
  \qquad\text{或等价地}\qquad
  Y\supset X_n\quad\text{对每个 \(n\)}.
\]
的情形. 此时须证明 \(T=S\), 即 \(Y=X\).

由对每个 \(n\) 有 \(Y\supset X_n\), 再结合
\(X\) 局部 Noether 这一事实, 可知在
\(\epsilon(S)\) 的每一点的一个邻域内, \(Y\) 为 \(X\) 的一个诱导开子概形.%
\footnote{\label{note:sga3-VIB-64}编者注: 例如见
\hyperref[cor:sga3-VIB-6.2.6]{6.2.6} 的证明.}
因此存在 \(X\) 的一个包含 \(\epsilon(S)\) 的诱导开子概形 \(U\),
使得 \(U\subset Y\). 由
EGA~IV\(_3\), 11.10.10, 由于 \(X/S\) 的纤维整,
\(U\) 在 \(X\) 中概形稠密. 由于 \(Y\) 为包含 \(U\) 的闭
子概形, 得到 \(Y=X\). 这就在 \(S\) 为 Noether 概形时证明了
\hyperref[cor:sga3-VIB-6.5.3]{6.5.3}, 从而也证明了
\hyperref[thm:sga3-VIB-6.5.2]{6.5.2}.

一般情形可化归为上述情形. 对每个
\(s\in S\), 存在 \(s\) 的仿射开邻域 \(U\)
和 \(\epsilon(s)\) 的仿射开邻域 \(V\),
使得 \(f(V)\subset U\). 则 \(f(V)\) 是包含于 \(U\) 的
\(s\) 的开邻域. 设 \(S'\) 为 \(s\) 的一个仿射
开邻域, 包含于
\[
  \epsilon^{-1}(V)\cap f(V),
\]
并令
\[
  X'=V\cap f^{-1}(S').
\]
则 \(X'\) 和 \(S'\) 分别为 \(X\) 和 \(S\) 的仿射开子集,
且 \(X'/S'\) 有一个截面. 由
\SGARefLink{thm:sga3-VIB-6.5.2}{6.5.2} 和 \SGARefLink{cor:sga3-VIB-6.5.3}{6.5.3} 的局部性, 可以假设 \(S=S'\).

因此 \(X'\) 为 \(X\) 的一个包含
\(\epsilon(S)\) 的仿射开子集. 由于每个纤维 \(X_s\) 均不可约,
\(X'_s\) 在 \(X_s\) 中概形稠密. 因此, 由
EGA~IV\(_3\), 11.10.10, \(X'\) 在 \(X\) 中概形稠密;
同样, 对每个基变换 \(S_1\to S\),
\[
  X'\times_S S_1
\]
在 \(X\times_S S_1\) 中概形稠密.

因此有
\phantomsection\label{eq:sga3-VIB-6.5.3-open-restriction}%
\[
  \prod_{X/S}Y/X
  =
  \prod_{X'/S}Y'/X',
  \qquad Y'=Y\cap X'.
\]
这将我们化归为 \(X=X'\) 的情形, 因而可以假设
\(S\) 和 \(X\) 都仿射. 此外, 若
\[
  X=\operatorname{Spec}B
\]
且 \(J\) 为 \(B\) 中定义 \(Y\) 的理想, 则 \(J\) 是其有限生成
子理想的滤过余极限. 因此 \(Y\)
是 \(X\) 的一族在 \(X\) 上有限呈示的闭子概形 \(Y_i\)
的交, 从而
\phantomsection\label{eq:sga3-VIB-6.5.3-finite-presentation-reduction}%
\[
  \prod_{X/S}Y/X
  =
  \bigcap_i\prod_{X/S}Y_i/X.
\]
因此, 为证明 \hyperref[thm:sga3-VIB-6.5.2]{6.5.2},
只需考虑 \(Y\) 在 \(X\) 上有限呈示,
且 \(S\) 和 \(X\) 仿射的情形.

此时 \(S\) 上的 \(X\) 和 \(Y\) 可由一个
Noether 概形 \(S_0\) 上的类似情形 \(X_0,Y_0\)
经基变换 \(S\to S_0\) 得到. 这将我们化归为已经处理的
Noether 情形, 从而完成了
\hyperref[thm:sga3-VIB-6.5.2]{6.5.2} 和
\hyperref[cor:sga3-VIB-6.5.3]{6.5.3} 的证明.

% 翻译统计: 源行数 190; 输出行数 188 (含本统计行); 行内/行间数学片段 134; 翻译占位 0. 底本疑点另报主线程.

% 独立复核检查点：已逐段对照英文源文，覆盖正文、假设、证明、脚注与源注。
% 行内/陈列数学及全部引用键机械对照通过；仅统一群的invariant为“正规”，保留源文疑点另报。
